% Volume VIII: Applications and Case Structures
% Part XV. Domains of Difficult Judgment
% Chapter 88: Employment and Institutional Reputation
% Spine: proof-spines/ch088-spine.md

\chapter{Employment and Institutional Reputation}
\label{ch:ch088}

Workplace investigations often proceed under incomplete evidence, uneven power, and intense pressure for rapid resolution. This chapter examines employment, academic, and professional settings in which allegations of harassment, retaliation, discrimination, incompetence, or disloyalty are filtered through managerial risk, confidentiality rules, and institutional reputation. The key distinction is between protective interim action and premature final judgment.

That distinction matters because organizations often need to reassign duties, separate parties, or pause access before they can know enough to impose lasting sanction. The chapter asks when such interim measures are justified, when reputational or career consequences outrun the evidence, and how institutions can avoid turning uncertainty into a self-protective narrative that quietly settles the case before any genuine testing occurs.
